\section{Результаты моделирования}

\subsection{Параметры модели}

Численное моделирование проводится для радиального подшипника скольжения геометрии~\varGeomLetter{} ($R = \varRmm$~мм, $c = \varCmkm$~мкм, $L = \varLmm$~мм) с \varDimpleTypeGenitive{} микрорельефом. Исходные данные расчёта сведены в таблицу~\ref{tab:input_data}.

\begin{table}[H]
\small
\caption{Исходные данные для расчёта}
\label{tab:input_data}
\begin{tabularx}{\textwidth}{|X|c|c|}
\hline
\centering\textbf{Параметр} & \textbf{Значение} & \textbf{Размерность} \\
\hline
Радиус вала $R$ & \varRmm & мм \\
\hline
Радиальный зазор $c$ & \varCmkm & мкм \\
\hline
Длина подшипника $L$ & \varLmm & мм \\
\hline
Относительная длина $L/D$ & \varLD & -- \\
\hline
Полуось по оси $Z$, $a$ & \varDimpleAmm & мм \\
\hline
Полуось по оси $\varphi$, $b$ & \varDimpleBmm & мм \\
\hline
Глубина углубления $h_p$ & \varDimpleHpmkm & мкм \\
\hline
Относительная глубина $h_p/c$ & \varDimpleHpRatio & -- \\
\hline
Число углублений по окружности $N_\varphi$ & \varNphi & шт. \\
\hline
Число рядов по длине $N_Z$ & \varNZ & шт. \\
\hline
Область текстурирования & \varPhiRange & -- \\
\hline
Частота вращения $n$ & \varN & об/мин \\
\hline
Вязкость смазки $\eta$ & \varEta & Па$\cdot$с \\
\hline
Диапазон эксцентриситета $\varepsilon$ & 0{,}05-0{,}80 & -- \\
\hline
Размер сетки & $500 \times 500$ & -- \\
\hline
Параметр релаксации $\omega$ & 1{,}5 & -- \\
\hline
Критерий сходимости $\delta$ & $< 10^{-5}$ & -- \\
\hline
Макс. число итераций & 50000 & -- \\
\hline
\end{tabularx}
\end{table}

При заданной частоте вращения линейная скорость поверхности вала равна $U = \pi D n / 60 = \pi \cdot 0{,}070 \cdot 2980 / 60 \approx 10{,}93$~м/с. Моделирование выполнено для 15 значений эксцентриситета в указанном диапазоне.

\subsection{Распределение зазора и давления}

На рисунке~\ref{fig:pressure_2d} показано распределение безразмерного давления $P(\varphi)$ в среднем сечении ($Z = 0$) при $\varepsilon = 0{,}6$ для гладкого и текстурированного подшипников.

\begin{figure}[H]
\centering
\includegraphics[width=\textwidth]{figures/pressure_2d.png}
\caption{Распределение давления $P(\varphi)$ при $Z=0$ для гладкого и текстурированного подшипников ($\varepsilon = 0{,}6$)}
\label{fig:pressure_2d}
\end{figure}

Кривая давления гладкого подшипника имеет классический вид с одним выраженным максимумом вблизи зоны минимального зазора ($\varphi \approx \pi$). В сужающейся части зазора давление постепенно нарастает, а в расширяющейся -- обращается в нуль в соответствии с условием кавитации. Зона положительного давления занимает приблизительно половину окружности подшипника. Введение микрорельефа сферического типа порождает локальные осцилляции давления, приуроченные к каждому углублению. Профиль сферической шапки, аппроксимирующий сферический сегмент, обеспечивает плавные переходы давления, аналогичные эллипсоидальному типу, но с иным распределением кривизны. На входной кромке каждого углубления формируется локальный пик, а на выходной -- понижение давления. При этом суммарный уровень давления текстурированного подшипника превышает уровень гладкого, что обусловливает прирост несущей способности.

Профиль зазора $H(\varphi)$ в среднем сечении при $\varepsilon = 0{,}6$ приведён на рисунке~\ref{fig:clearance_2d}.

\begin{figure}[H]
\centering
\includegraphics[width=\textwidth]{figures/clearance_2d.png}
\caption{Распределение зазора $H(\varphi)$ при $Z=0$ для подшипника с \varDimpleTypeGenitive{} углублениями ($\varepsilon = 0{,}6$)}
\label{fig:clearance_2d}
\end{figure}

Плавное изменение зазора по закону $H = 1 + \varepsilon\cos\varphi$ дополняется локальными увеличениями в зонах расположения углублений сферического профиля. Максимальная безразмерная глубина каждого углубления составляет $\Delta H_{\max} = h_p/c = 0{,}2$.

Трёхмерные распределения зазора и давления при $\varepsilon = 0{,}6$ представлены на рисунках~\ref{fig:clearance_3d} и~\ref{fig:pressure_3d}.

\begin{figure}[H]
\centering
\includegraphics[width=\textwidth]{figures/clearance_3d.png}
\caption{Трёхмерное распределение зазора $H(\varphi, Z)$ ($\varepsilon = 0{,}6$)}
\label{fig:clearance_3d}
\end{figure}

Рисунок~\ref{fig:clearance_3d} демонстрирует упорядоченную структуру микрорельефа: 8 углублений по окружности и 11 рядов по длине подшипника в области $90^\circ - 270^\circ$.

\begin{figure}[H]
\centering
\includegraphics[width=\textwidth]{figures/pressure_3d.png}
\caption{Трёхмерное распределение давления $P(\varphi, Z)$ ($\varepsilon = 0{,}6$)}
\label{fig:pressure_3d}
\end{figure}

Рисунок~\ref{fig:pressure_3d} показывает, что наряду с основным максимумом давления, формируемым клиновым зазором макрогеометрии, присутствуют локальные возмущения на каждом углублении. Сферическая форма профиля создаёт на входной кромке углубления сужение зазора, формирующее дополнительный рост давления -- микрогидродинамический эффект. На выходной кромке происходит расширение зазора и снижение давления, однако кавитационное ограничение ($P \geq 0$) обеспечивает положительный суммарный вклад каждого углубления.

Из анализа трёхмерных распределений следует, что распределение давления по длине подшипника существенно неоднородно: максимальные значения достигаются в среднем сечении ($Z = 0$), тогда как у торцов ($Z = \pm 1$) давление стремится к нулю из-за утечек. Воздействие текстурирования проявляется сильнее всего в центральной части, где фоновое давление наибольшее. Углубления, расположенные вблизи торцов подшипника, дают менее выраженный эффект вследствие низкого уровня базового давления в этих зонах.

\subsection{Сравнительный анализ интегральных показателей}

В таблицах~\ref{tab:smooth_results} и~\ref{tab:textured_results} представлены интегральные характеристики гладкого подшипника и подшипника с углублениями сферического профиля для выбранных значений эксцентриситета.

\begin{table}[H]
\small
\centering
\caption{Характеристики гладкого подшипника}
\label{tab:smooth_results}
\begin{tabular}{|c|c|c|c|}
\hline
$\boldsymbol{\varepsilon}$ & $\boldsymbol{F}$\textbf{, Н} & $\boldsymbol{\mu}$ & $\boldsymbol{Q}$\textbf{, л/с} \\
\hline
0{,}05 & 259{,}5 & 0{,}2291 & 0{,}2398 \\
\hline
0{,}16 & 835{,}7 & 0{,}0722 & 0{,}2395 \\
\hline
0{,}26 & 1478{,}5 & 0{,}0420 & 0{,}2389 \\
\hline
0{,}37 & 2261{,}3 & 0{,}0287 & 0{,}2380 \\
\hline
0{,}48 & 3286{,}8 & 0{,}0211 & 0{,}2370 \\
\hline
0{,}59 & 4764{,}4 & 0{,}0160 & 0{,}2360 \\
\hline
0{,}69 & 7222{,}5 & 0{,}0120 & 0{,}2351 \\
\hline
0{,}80 & 12218{,}3 & 0{,}0087 & 0{,}2344 \\
\hline
\end{tabular}
\end{table}

\begin{table}[H]
\small
\centering
\caption{Характеристики подшипника с \varDimpleTypeGenitive{} углублениями}
\label{tab:textured_results}
\begin{tabular}{|c|c|c|c|}
\hline
$\boldsymbol{\varepsilon}$ & $\boldsymbol{F}$\textbf{, Н} & $\boldsymbol{\mu}$ & $\boldsymbol{Q}$\textbf{, л/с} \\
\hline
0{,}05 & 346{,}5 & 0{,}1698 & 0{,}2429 \\
\hline
0{,}16 & 912{,}9 & 0{,}0653 & 0{,}2426 \\
\hline
0{,}26 & 1577{,}0 & 0{,}0388 & 0{,}2420 \\
\hline
0{,}37 & 2406{,}0 & 0{,}0266 & 0{,}2411 \\
\hline
0{,}48 & 3515{,}8 & 0{,}0194 & 0{,}2401 \\
\hline
0{,}59 & 5233{,}8 & 0{,}0143 & 0{,}2390 \\
\hline
0{,}69 & 8451{,}1 & 0{,}0101 & 0{,}2380 \\
\hline
0{,}80 & 17046{,}7 & 0{,}0061 & 0{,}2370 \\
\hline
\end{tabular}
\end{table}

С увеличением эксцентриситета нагрузочная способность $F$ нелинейно возрастает для обоих вариантов подшипника. Во всём диапазоне $\varepsilon$ текстурированный подшипник демонстрирует более высокие значения нагрузочной способности. При больших эксцентриситетах ($\varepsilon = 0{,}80$) прирост особенно значителен: нагрузочная способность текстурированного подшипника (17046{,}7~Н) превосходит значение для гладкого (12218{,}3~Н) на 39{,}5\%.

Коэффициент трения $\mu$ уменьшается с ростом $\varepsilon$ и во всех точках ниже у текстурированного подшипника. Расход смазки $Q$ демонстрирует слабую зависимость от эксцентриситета и несколько выше для текстурированного варианта, что объясняется увеличением эффективного объёма смазочного слоя в области углублений.

Количественная оценка влияния текстурирования при $\varepsilon = 0{,}6$ приведена в таблице~\ref{tab:comparison}.

\begin{table}[H]
\small
\centering
\caption{Сравнение характеристик при $\varepsilon = 0{,}6$}
\label{tab:comparison}
\begin{tabular}{|l|c|c|c|}
\hline
\textbf{Характеристика} & \textbf{Гладкий} & \textbf{С углублениями} & $\boldsymbol{\Delta}$\textbf{, \%} \\
\hline
Нагрузка $F$, Н & 5028{,}5 & 5552{,}5 & $+10{,}42$ \\
\hline
Коэфф. трения $\mu$ & 0{,}01538 & 0{,}01365 & $-11{,}22$ \\
\hline
Расход $Q$, л/с & 0{,}2359 & 0{,}2389 & $+1{,}27$ \\
\hline
\end{tabular}
\end{table}

Применение углублений сферического профиля позволяет повысить несущую способность подшипника на 10{,}42\% благодаря микрогидродинамическому эффекту на кромках углублений. Снижение коэффициента трения на 11{,}22\% обусловлено ростом средней толщины смазочной плёнки в зоне текстурирования. Увеличение расхода смазки составляет 1{,}27\%, что благоприятно для теплоотвода и не предъявляет чрезмерных требований к системе маслоснабжения.

Зависимости интегральных характеристик от эксцентриситета для обоих вариантов подшипника представлены на рисунках~\ref{fig:load_vs_eps}-\ref{fig:flow_vs_eps}.

\begin{figure}[H]
\centering
\includegraphics[width=\textwidth]{figures/load_vs_epsilon.png}
\caption{Нагрузочная способность $F(\varepsilon)$ для гладкого и текстурированного подшипников}
\label{fig:load_vs_eps}
\end{figure}

Во всём диапазоне эксцентриситетов кривая $F(\varepsilon)$ текстурированного подшипника расположена выше кривой гладкого. Различие между кривыми нарастает при увеличении $\varepsilon$: при $\varepsilon = 0{,}05$ разность составляет около 87~Н, тогда как при $\varepsilon = 0{,}80$ она достигает порядка 4800~Н. Такое поведение определяется степенной зависимостью давления от зазора ($P \sim h^{-3}$), усиливающей влияние текстурирования при уменьшении толщины смазочного слоя.

\begin{figure}[H]
\centering
\includegraphics[width=\textwidth]{figures/friction_vs_epsilon.png}
\caption{Коэффициент трения $\mu(\varepsilon)$ для гладкого и текстурированного подшипников}
\label{fig:friction_vs_eps}
\end{figure}

При всех значениях $\varepsilon$ коэффициент трения текстурированного подшипника остаётся ниже, чем у гладкого. Обе зависимости монотонно убывают, стремясь к нулю при $\varepsilon \to 1$. Максимальная абсолютная разница наблюдается при малых эксцентриситетах: при $\varepsilon = 0{,}05$ она составляет $\Delta\mu = 0{,}059$, а при $\varepsilon = 0{,}80$ снижается до $\Delta\mu = 0{,}003$. В относительных единицах снижение трения остаётся сопоставимым по всему диапазону.

\begin{figure}[H]
\centering
\includegraphics[width=\textwidth]{figures/flow_vs_epsilon.png}
\caption{Расход смазки $Q(\varepsilon)$ для гладкого и текстурированного подшипников}
\label{fig:flow_vs_eps}
\end{figure}

Расход смазки текстурированного подшипника превышает расход гладкого на величину порядка 1-1{,}3\%, причём эта разница практически постоянна по всему диапазону эксцентриситетов.

Таким образом, использование сферического микрорельефа обеспечивает одновременный рост нагрузочной способности и уменьшение коэффициента трения при малом увеличении расхода смазки. Данный эффект усиливается с ростом эксцентриситета, что представляет особый интерес для тяжелонагруженных узлов.

\subsection*{Выводы по главе 3}
\phantomsection
\addcontentsline{toc}{subsection}{Выводы по главе 3}

Проведён расчёт характеристик радиального подшипника скольжения со сферическим микрорельефом для 15 значений эксцентриситета. Определены поля давления и зазора, построены графические зависимости нагрузочной способности, коэффициента трения и расхода смазки от эксцентриситета. Сравнение характеристик при $\varepsilon = 0{,}6$ выявило прирост нагрузочной способности на $10{,}42$\%, снижение коэффициента трения на $11{,}22$\% и увеличение расхода на $1{,}27$\%, что свидетельствует об эффективности текстурирования углублениями сферического типа.
